% ================= 4. CÀI ĐẶT - 2 DESIGN PATTERN =================
\section{Cài đặt -- 02 Design Pattern}

\begin{frame}{Kiến trúc cài đặt và các mẫu thiết kế trong code}
\begin{columns}[T]
\column{0.4\textwidth}
\centering
\begin{tikzpicture}[
  a/.style={-{Stealth[length=2mm]}, draw=black!65, thick}]
\node[pk=navy] (ui) at (0,4.2) {\textbf{HotelCRM.WinForms} \ {\tiny(giao diện)}\\{\tiny Form, Program.cs, ErrorHandler}};
\node[pk=tvA] (bl) at (0,2.85) {\textbf{HotelCRM.Business} \ {\tiny(nghiệp vụ)}\\{\tiny Services, Queries, Strategy, Factory}};
\node[pk=tvB] (dm) at (0,1.5) {\textbf{HotelCRM.Domain} \ {\tiny(lõi)}\\{\tiny thực thể, State, quy tắc, interface}};
\node[pk=common] (da) at (0,0.15) {\textbf{HotelCRM.Data} \ {\tiny(dữ liệu)}\\{\tiny JsonRepository, Record, Mapper}};
\draw[a] (ui) -- (bl); \draw[a] (bl) -- (dm); \draw[a] (da) -- (dm);
\draw[a,dashed] (ui.east) -- ++(0.35,0) |- node[pos=0.25,right,font=\tiny,align=left]{chỉ\\\code{Program.cs}} (da.east);
\end{tikzpicture}

\vspace{1mm}
{\tiny Domain không tham chiếu project nào $\Rightarrow$ không có SQL/UI trong lõi. \code{HotelCRM.Tests} tham chiếu Domain, Business, Data và chạy không cần giao diện.\par}
\column{0.58\textwidth}
\tiny
\begin{tabularx}{\linewidth}{@{}L{1.9cm}L{0.75cm}YL{1.3cm}@{}}
\toprule
\textbf{Mẫu (nhóm GoF)} & \textbf{TV} & \textbf{Lớp tham gia} & \textbf{Trạng thái} \\
\midrule
\rowcolor{tvA!10} \textbf{State} (hành vi) & TV1 & \code{SupportTicket}, \code{TicketState} + 4 state & \stOK \\
Factory Method (khởi tạo) & TV1 & \code{TicketFactory} + 3 factory & \stOK \\
\rowcolor{tvB!10} \textbf{Strategy} (hành vi) & TV2 & \code{IAssignmentStrategy} + 3 strategy; Context \code{TicketService} (TV1) & \stOK \ cầu nối \\
Observer (hành vi) & TV1/TV2 & Subject \code{SupportTicket}; 3 observer & Subject \stOK; observer: GĐ 4 \\
Singleton (khởi tạo) & Chung & \code{AppConfig.Instance} & \stOK \\
\bottomrule
\end{tabularx}

\vspace{2mm}
\begin{block}{Hai mẫu trình bày chi tiết}
\scriptsize
\textbf{State} (TV1): vòng đời ticket. \textbf{Strategy} (TV2): thuật toán phân công -- là \textbf{cầu nối} vì Context nằm ở phân hệ TV1, các ConcreteStrategy nằm ở phân hệ TV2.
\end{block}
\end{columns}
\end{frame}

% ---------- STATE ----------
\begin{frame}{Pattern 1 -- State (TV1): bài toán và sơ đồ trạng thái ticket}
\begin{columns}[c]
\column{0.42\textwidth}
\scriptsize
\textbf{Bài toán.} Ticket có 4 trạng thái; mỗi thao tác (Start, Resolve, Close, Reopen) hợp lệ ở trạng thái này nhưng bị cấm ở trạng thái khác (quy tắc 1, 2).\\[1mm]
\textbf{Nếu không dùng mẫu:} mỗi hàm của \code{SupportTicket} là một chuỗi \code{if/switch} theo trạng thái; thêm trạng thái mới phải sửa cả 4 hàm.\\[1mm]
\textbf{Giải pháp:} mỗi trạng thái là một lớp. \code{SupportTicket} ủy quyền thao tác cho đối tượng \code{State} hiện tại; lớp trạng thái tự quyết định chuyển sang đâu hoặc ném \code{BusinessRuleException}.
\column{0.56\textwidth}
\centering
\begin{tikzpicture}[
  s/.style={draw=tvA, fill=tvA!10, rounded corners=6pt, minimum width=2.2cm, minimum height=0.8cm, font=\scriptsize\bfseries},
  t/.style={-{Stealth[length=2mm]}, draw=black!70, font=\tiny}]
\node[circle, fill=black, inner sep=2.5pt] (i) at (-1.6,2.2) {};
\node[s] (o) at (0,2.2) {Open};
\node[s] (p) at (3.6,2.2) {InProgress};
\node[s] (r) at (3.6,0) {Resolved};
\node[s] (c) at (0,0) {Closed};
\node[circle, draw=black, inner sep=1.5pt] (e) at (-1.6,0) {}; \fill (e) circle (1.6pt);
\draw[t] (i) -- node[above]{CreateTicket} (o);
\draw[t] (o) -- node[above]{Start [có Agent]} node[below]{Agent được giao} (p);
\draw[t] (p) -- node[right,align=left]{Resolve\\{}[ghi chú không rỗng]} (r);
\draw[t] (r.west) -- node[above]{Close} node[below]{Manager} (c.east);
\draw[t] (r.north west) to[bend left=35] node[left,align=right,pos=0.45]{Reopen [now $\le$ ResolvedAt\\+ ReopenDays] -- Manager} (p.south west);
\draw[t] (c) -- (e);
\draw[t] (o.north) to[out=110,in=70,looseness=6] node[above]{AssignTo (Manager)} (o.north east);
\draw[t] (p.north) to[out=110,in=70,looseness=6] node[above]{AssignTo (gán lại)} (p.north east);
\end{tikzpicture}

\vspace{1mm}
{\tiny Mọi mũi tên không vẽ (ví dụ Resolve khi Closed) $\Rightarrow$ lớp cha \code{TicketState} ném \code{BusinessRuleException}.\par}
\end{columns}
\end{frame}

\begin{frame}{Pattern 1 -- State: vai trò các lớp}
\begin{columns}[c]
\column{0.5\textwidth}
\centering
\resizebox{\linewidth}{!}{%
\begin{tikzpicture}[every node/.append style={anchor=north}]
\path (0,4) node[umlcls=tvA] (ctx) {{\centering\textbf{SupportTicket} \textcolor{warn}{«Context»}\par}\nodepart{second}-State: TicketState\nodepart{third}+Start(now)\\+Resolve(note, now)\\+Close(now)\\+Reopen(now, reopenDays)\\\textasciitilde MoveTo(next, now)};
\path (5.0,4) node[umlcls=tvA] (st) {{\centering\textbf{\textit{TicketState}} \textcolor{warn}{«State»}\par}\nodepart{second}+Name, IsActive, IsClosed\nodepart{third}\textasciitilde Start(t, now)\\\textasciitilde Resolve(t, note, now)\\\textasciitilde Close(t, now)\\\textasciitilde Reopen(t, now, days)};
\foreach \x/\n/\m in {2.4/Open/Start, 4.3/InProgress/Resolve, 6.2/Resolved/{Close, Reopen}, 8.1/Closed/--}{
  \path (\x,0.6) node[umlcls=tvA,minimum width=1.75cm] (z\n) {{\centering\textbf{\n State}\par}\nodepart{second}~\nodepart{third}\m};
  \draw[thick,black!70] (z\n.north) -- ++(0,0.3);}
\draw[thick,black!70] (2.4,0.9) -- (8.1,0.9);
\draw[inh] (5.0,0.9) -- (st.south);
\draw[comp] (ctx.east |- 0,3.5) -- node[mult,above,pos=0.8]{1} (st.west |- 0,3.5);
\node[font=\tiny, text=warn] at (6.2,-1.0) {«ConcreteState»};
\end{tikzpicture}}
\column{0.48\textwidth}
\scriptsize
\begin{tabularx}{\linewidth}{@{}L{1.7cm}Y@{}}
\toprule
\textbf{Vai trò GoF} & \textbf{Lớp trong code} \\
\midrule
Context & \code{SupportTicket} -- giữ \code{State}, ủy quyền 4 thao tác \\
State & \code{TicketState} (abstract) -- mặc định \textbf{từ chối} mọi thao tác \\
ConcreteState & \code{OpenState}, \code{InProgressState}, \code{ResolvedState}, \code{ClosedState} -- chỉ override thao tác hợp lệ \\
Chuyển trạng thái & ConcreteState gọi \code{ticket.MoveTo(new XxxState(), now)} \\
Khôi phục từ file & \code{TicketState.FromName(name)} \\
\bottomrule
\end{tabularx}

\vspace{2mm}
\textbf{Cơ chế phối hợp:} đa hình (gọi \code{State.Start()} không biết lớp cụ thể) + ủy quyền (Context chuyển việc cho State) + State tạo State kế tiếp.
\end{columns}
\end{frame}

\begin{frame}[fragile]{Pattern 1 -- State: cài đặt trong code}
\lstset{basicstyle=\ttfamily\tiny}
\begin{columns}[T]
\column{0.5\textwidth}
\textbf{\scriptsize Context -- \code{Domain/Tickets/SupportTicket.cs}}
\begin{lstlisting}
public void Start(DateTime now)
{
    State.Start(this, now);   // delegate to the state
}

internal void MoveTo(TicketState next, DateTime now)
{
    TicketState previous = State;
    State = next;
    UpdatedAt = now;
    Notify(previous, next);   // Observer hook
}
\end{lstlisting}
\textbf{\scriptsize State -- \code{States/TicketState.cs}}
\begin{lstlisting}
internal virtual void Start(SupportTicket ticket, DateTime now)
{
    throw new BusinessRuleException(
        $"A ticket that is {Name} cannot be started.");
}
\end{lstlisting}
\column{0.48\textwidth}
\textbf{\scriptsize ConcreteState -- \code{States/OpenState.cs}}
\begin{lstlisting}
internal override void Start(SupportTicket ticket, DateTime now)
{
    if (ticket.AssignedAgent is null)
    {
        throw new BusinessRuleException(
            "The ticket must be assigned to an agent "
            + "before it can be started.");
    }
    ticket.MoveTo(new InProgressState(), now);
}
\end{lstlisting}
\textbf{\scriptsize ConcreteState -- \code{States/ResolvedState.cs}}
\begin{lstlisting}
internal override void Reopen(SupportTicket ticket,
    DateTime now, int reopenDays)
{
    DateTime resolvedAt = ticket.ResolvedAt ?? ticket.UpdatedAt;
    DateTime deadline = resolvedAt.AddDays(reopenDays);
    if (now > deadline) { throw new BusinessRuleException(...); }
    ticket.MoveTo(new InProgressState(), now);
}
\end{lstlisting}
\end{columns}
\end{frame}

\begin{frame}{Pattern 1 -- State: đánh giá và bài mở rộng}
\begin{columns}[T]
\column{0.5\textwidth}
\begin{block}{Hiệu quả}
\scriptsize
\begin{itemize}
  \item Quy tắc 1, 2 nằm đúng một chỗ cho mỗi trạng thái; \code{SupportTicket} không có \code{switch} theo trạng thái.
  \item Trạng thái cấm bị chặn \textbf{trong Domain}, kể cả khi gọi Service không qua giao diện.
  \item 6 test \code{TicketStateTests} kiểm tra cả chuyển hợp lệ và chuyển cấm.
\end{itemize}
\end{block}
\begin{block}{Nhược điểm / đánh đổi}
\scriptsize
\begin{itemize}
  \item Thêm 5 lớp cho 4 trạng thái; người đọc phải mở nhiều file để thấy toàn bộ máy trạng thái.
  \item Lưu file chỉ lưu tên trạng thái, cần bảng \code{FromName} để dựng lại.
\end{itemize}
\end{block}
\column{0.48\textwidth}
\begin{block}{Bài mở rộng: thêm \code{OnHoldState} (chờ khách phản hồi)}
\scriptsize
\begin{enumerate}
  \item Tạo lớp \code{OnHoldState : TicketState}, override \code{Start()} để quay lại InProgress.
  \item Thêm thao tác \code{Hold()} vào \code{TicketState} (mặc định ném lỗi) và override trong \code{InProgressState}.
  \item Thêm \textbf{1 dòng} vào \code{TicketState.FromName()} để nạp lại từ file.
\end{enumerate}
Các ConcreteState cũ không phải sửa; đúng tinh thần ``thêm lớp mới + đăng ký 1 dòng'' (OCP).
\end{block}
\end{columns}
\end{frame}

% ---------- STRATEGY ----------
\begin{frame}{Pattern 2 -- Strategy (TV2): bài toán và vai trò các lớp}
\begin{columns}[c]
\column{0.4\textwidth}
\scriptsize
\textbf{Bài toán.} Manager bấm \emph{Auto Assign}: hệ thống chọn Agent cho ticket. Có nhiều cách chọn và khách sạn muốn đổi cách chọn \textbf{không cần sửa code}:
\begin{itemize}
  \item \code{LeastLoaded}: Agent ít ticket đang mở nhất
  \item \code{DepartmentMatch}: ưu tiên Agent thuộc bộ phận xử lý của loại ticket (\code{GUEST}/\code{HOUSE}/\code{ENG}), không còn ai thì chọn trong toàn bộ
  \item \code{RoundRobin}: xoay vòng theo Id
\end{itemize}
\textbf{Cầu nối 2 phân hệ:} Context \code{TicketService} (TV1) chỉ biết interface; 3 lớp cụ thể do TV2 viết.
\column{0.58\textwidth}
\centering
\resizebox{\linewidth}{!}{%
\begin{tikzpicture}[every node/.append style={anchor=north}]
\path (0,4.2) node[umlcls=tvA] (ctx) {{\centering\textbf{TicketService} \textcolor{warn}{«Context»}\par}\nodepart{second}-\_strategy: IAssignmentStrategy\nodepart{third}+AutoAssign(actor, ticketId): Agent};
\path (6.0,4.2) node[umlcls=navy] (ist) {{\centering\textit{«interface»} \textbf{IAssignmentStrategy} \textcolor{warn}{«Strategy»}\par}\nodepart{second}~\nodepart{third}+Pick(departmentCode, candidates): Agent?};
\path (3.0,1.4) node[umlcls=tvB] (s1) {{\centering\textbf{LeastLoadedStrategy}\par}\nodepart{second}~\nodepart{third}+Pick(...)};
\path (6.0,1.4) node[umlcls=tvB] (s2) {{\centering\textbf{DepartmentMatchStrategy}\par}\nodepart{second}-\_leastLoaded\nodepart{third}+Pick(...)};
\path (9.0,1.4) node[umlcls=tvB] (s3) {{\centering\textbf{RoundRobinStrategy}\par}\nodepart{second}-\_lastPickedId: int\nodepart{third}+Pick(...)};
\path (0,1.4) node[umlcls=common] (wl) {{\centering\textit{«record»} \textbf{AgentWorkload}\par}\nodepart{second}+Agent, ActiveCount\nodepart{third}~};
\path (9.0,-0.6) node[draw=navy, fill=navy!8, font=\scriptsize, align=center] (pr) {\code{Program.cs}: bảng \code{tên -> strategy}\\chọn theo \code{appsettings.json}};
\draw[aggr] (ctx.east |- 0,3.6) -- node[mult,pos=0.85,above]{1} (ist.west |- 0,3.6);
\coordinate (j) at (6.0,1.75);
\draw[impl] (j) -- (ist.south);
\foreach \s in {s1,s2,s3}{\draw[dashed,black!70] (\s.north) |- (j);}
\draw[comp] (s2.south) -- ++(0,-0.4) -| node[mult,pos=0.3,below]{dùng lại} (s1.south);
\draw[dep] (ctx.south) -- node[mult,right]{tạo} (wl.north);
\draw[dep] (pr.north) -- (s3.south);
\node[font=\tiny, text=warn] at (6.0,-0.3) {«ConcreteStrategy» (TV2)};
\end{tikzpicture}}
\end{columns}
\end{frame}

\begin{frame}[fragile]{Pattern 2 -- Strategy: cài đặt trong code}
\lstset{basicstyle=\ttfamily\tiny}
\begin{columns}[T]
\column{0.53\textwidth}
\textbf{\scriptsize Context -- \code{Business/Services/TicketService.cs} (TV1)}
\begin{lstlisting}
public Agent AutoAssign(Employee actor, int ticketId)
{
    actor.EnsureCanAssign();                  // rule 11
    SupportTicket ticket = _tickets.GetRequired(ticketId);
    List<AgentWorkload> candidates = new();
    foreach (Agent agent in GetAgents())
    {
        int activeCount = CountActiveTickets(agent.Id, ticket.Id);
        if (agent.CanTake(activeCount, _settings.MaxActiveTickets))
            candidates.Add(new AgentWorkload(agent, activeCount));
    }
    Agent? chosen = _strategy.Pick(ticket.HandlingDepartmentCode,
                                   candidates);   // polymorphic call
    ...
    ticket.AssignTo(chosen, chosenCount, _settings.MaxActiveTickets);
    _tickets.Save(ticket);
    return ticket.AssignedAgent!;
}
\end{lstlisting}
\column{0.45\textwidth}
\textbf{\scriptsize ConcreteStrategy -- \code{Business/Assignment/} (TV2)}
\begin{lstlisting}
public sealed class DepartmentMatchStrategy
    : IAssignmentStrategy
{
    private readonly LeastLoadedStrategy
        _leastLoaded = new();

    public Agent? Pick(string departmentCode,
        IReadOnlyList<AgentWorkload> candidates)
    {
        List<AgentWorkload> inDepartment = candidates
            .Where(c => string.Equals(
                c.Agent.Department.Code, departmentCode,
                StringComparison.OrdinalIgnoreCase))
            .ToList();
        return _leastLoaded.Pick(departmentCode,
            inDepartment.Count > 0 ? inDepartment : candidates);
    }
}
\end{lstlisting}
\textbf{\scriptsize Chọn strategy -- \code{Program.cs} + \code{appsettings.json}}
\begin{lstlisting}
["LeastLoaded"]     = new LeastLoadedStrategy(),
["DepartmentMatch"] = new DepartmentMatchStrategy(),
["RoundRobin"]      = new RoundRobinStrategy(),
// appsettings.json: "DefaultAssignmentStrategy": "LeastLoaded"
\end{lstlisting}
\end{columns}
\end{frame}

\begin{frame}{Pattern 2 -- Strategy: luồng chạy Auto Assign (xuyên 2 phân hệ)}
\centering
\resizebox{\textwidth}{!}{%
\begin{tikzpicture}[
  msg/.style={-{Stealth[length=2mm]}, draw=black!75, font=\tiny},
  ret/.style={dashed,-{Stealth[length=1.6mm]}, draw=black!55, font=\tiny}]
\foreach \x/\n/\t/\c in {0/f/{:TicketListForm}/navy, 3.2/s/{:TicketService}/tvA, 6.4/g/{:IAssignmentStrategy}/tvB, 9.6/t/{:SupportTicket}/tvA, 12.8/a/{:Agent}/common, 16.0/r/{:IRepository}/common}{
  \node[ll=\c] (\n) at (\x,0) {\t};
  \draw[black!40, dashed] (\n.south) -- (\x,-6.0);}
\node[font=\scriptsize] (m) at (-2.2,0) {Manager};
\draw[msg] (-2.2,-0.9) -- node[above]{bấm Auto Assign} (0,-0.9);
\draw[msg] (0,-1.4) -- node[above]{\code{AutoAssign(actor, id)}} (3.2,-1.4);
\draw[msg] (3.2,-1.8) to[out=0,in=0,looseness=4] node[right,align=left]{\code{EnsureCanAssign()}, đếm tải,\\lọc \code{Agent.CanTake()}} (3.2,-2.4);
\draw[msg] (3.2,-2.9) -- node[above]{\code{Pick(dept, candidates)}} (6.4,-2.9);
\draw[ret] (6.4,-3.4) -- node[above]{\code{Agent?}} (3.2,-3.4);
\draw[msg] (3.2,-3.9) -- node[above]{\code{AssignTo(agent, active, max)}} (9.6,-3.9);
\draw[msg] (9.6,-4.4) -- node[above]{\code{EnsureCanTake()}} (12.8,-4.4);
\draw[msg] (3.2,-5.0) -- node[above,pos=0.8]{\code{Save(ticket)}} (16.0,-5.0);
\draw[ret] (3.2,-5.6) -- node[above]{Agent được chọn} (0,-5.6);
\node[draw=tvB, fill=tvB!8, font=\tiny, align=left, anchor=north west] at (5.4,-0.45) {lớp cụ thể do \code{Program.cs} chọn:\\\code{LeastLoaded} | \code{DepartmentMatch} | \code{RoundRobin}};
\end{tikzpicture}}

\vspace{1mm}
{\scriptsize Strategy chỉ \textbf{chọn}; quyết định cuối vẫn ở Domain: \code{SupportTicket.AssignTo()} kiểm tra trạng thái và \code{Agent.EnsureCanTake()} kiểm tra giới hạn tải (quy tắc 1, 3). Strategy trả \code{null} $\Rightarrow$ Domain ném ``No agent is available''.\par}
\end{frame}

\begin{frame}{Pattern 2 -- Strategy: đánh giá và bài mở rộng}
\begin{columns}[T]
\column{0.5\textwidth}
\begin{block}{Hiệu quả}
\scriptsize
\begin{itemize}
  \item Đổi thuật toán bằng 1 dòng trong \code{appsettings.json}, không build lại.
  \item TV1 viết và test \code{TicketService} với strategy giả; TV2 test strategy không cần code ticket (\code{AgentWorkload} chỉ chứa \code{Agent} và số đếm).
  \item Test \code{LeastLoaded_Pick_ReturnsAgentWithFewestActive}.
\end{itemize}
\end{block}
\begin{block}{Nhược điểm / đánh đổi}
\scriptsize
\begin{itemize}
  \item Thêm 1 interface + 3 lớp cho một hàm chọn; với 1 thuật toán duy nhất thì là thừa.
  \item \code{RoundRobinStrategy} có trạng thái (\code{_lastPickedId}) nên không dùng chung an toàn giữa nhiều luồng; mất khi tắt app.
\end{itemize}
\end{block}
\column{0.48\textwidth}
\begin{block}{Bài mở rộng: thêm \code{SkillBasedStrategy}}
\scriptsize
\begin{enumerate}
  \item Tạo lớp \code{SkillBasedStrategy : IAssignmentStrategy}, cài \code{Pick()} (ví dụ ưu tiên Agent có kỹ năng ngôn ngữ của khách, hòa thì ít tải nhất).
  \item Đăng ký \textbf{1 dòng} trong bảng của \code{Program.cs}: \code{["SkillBased"] = new SkillBasedStrategy()}.
  \item Đổi \code{DefaultAssignmentStrategy} thành \code{SkillBased}.
\end{enumerate}
\code{TicketService} và 3 strategy cũ không sửa dòng nào (OCP). Tên lạ trong cấu hình $\to$ báo lỗi lúc khởi động, liệt kê tên hợp lệ.
\end{block}
\end{columns}
\end{frame}
